%=====================================================================
%  Курсовая работа. Рациональная параметризация конических сечений
%  Булгакова М. В., МПГУ, кафедра геометрии им. Л. С. Атанасяна, 2026
%
%  pdfLaTeX
%=====================================================================
\documentclass[a4paper,14pt]{extreport}

\usepackage{cmap}
\usepackage[T2A]{fontenc}
\usepackage[utf8]{inputenc}
\usepackage[english,russian]{babel}

% \usepackage{tempora}

\usepackage{amsmath,amssymb,amsthm}
\usepackage{geometry}
\usepackage{setspace}
\usepackage{indentfirst}
\usepackage{enumitem}
\usepackage{graphicx}
\usepackage{caption}
\usepackage{tikz}
\usepackage{pgfplots}
\pgfplotsset{compat=1.18}
\usetikzlibrary{arrows.meta,positioning,calc}
\usepackage[hidelinks]{hyperref}
\usepackage{xurl}

%--------------------------- Оформление по ГОСТ ----------------------
\geometry{left=30mm,right=15mm,top=20mm,bottom=20mm}
\onehalfspacing
\setlength{\parindent}{1.25cm}
\captionsetup[figure]{labelsep=endash,name=Рисунок,font=normalsize}

% нумерация страниц снизу по центру
\pagestyle{plain}

% заголовки глав и разделов
\usepackage{titlesec}
\titleformat{\chapter}[block]
  {\normalfont\bfseries\filcenter}{\thechapter.}{0.6em}{}
\titlespacing*{\chapter}{0pt}{0pt}{16pt}
\titleformat{\section}[block]
  {\normalfont\bfseries\filcenter}{\thesection.}{0.5em}{}
\titlespacing*{\section}{0pt}{14pt}{10pt}

% теоремы: сквозная нумерация внутри главы
\theoremstyle{definition}
\newtheorem{thm}{Теорема}[chapter]
\newtheorem{lem}[thm]{Лемма}
\newtheorem{defn}[thm]{Определение}
\newtheorem{cor}[thm]{Следствие}
\newtheorem{prop}[thm]{Предложение}
\newtheorem{exmp}[thm]{Пример}
\newtheorem{exer}[thm]{Упражнение}
\renewcommand{\proofname}{Доказательство}

\newcommand{\answer}[1]{\par\medskip\noindent\textbf{Ответ:} #1}

\newcommand{\Q}{\mathbb{Q}}
\newcommand{\R}{\mathbb{R}}
\newcommand{\Z}{\mathbb{Z}}
\newcommand{\PP}{\mathbb{P}}
\providecommand{\gcdru}{}\renewcommand{\gcdru}{\mathop{\text{НОД}}\nolimits}

\begin{document}
\sloppy

%=====================================================================
%                         ТИТУЛЬНЫЙ ЛИСТ
%=====================================================================
\begin{titlepage}
\singlespacing
\centering
{\bfseries Министерство просвещения Российской Федерации}\\[2pt]
федеральное государственное бюджетное образовательное\\
учреждение высшего образования\\
<<Московский педагогический государственный университет>>\\[14pt]

Институт математики и информатики\\
Кафедра геометрии им.~Л.~С.~Атанасяна\\[26pt]

КУРСОВАЯ РАБОТА\\[14pt]

По дисциплине: \textbf{ГЕОМЕТРИЯ}\\[10pt]

На тему: \textbf{РАЦИОНАЛЬНАЯ ПАРАМЕТРИЗАЦИЯ}\\
\textbf{КОНИЧЕСКИХ СЕЧЕНИЙ}\\[14pt]

Код и направление подготовки: 44.03.05\\
Педагогическое образование (с двумя профилями подготовки)\\
Направленность (профиль) образовательной программы:\\
Математика и экономика\\[16pt]

\raggedleft
Выполнила:\\
Студентка 3 курса\\
группа №~305\\
Булгакова Мария Владимировна\\[14pt]

\raggedright
<<Допуск к защите>>\\[10pt]
\rule{6cm}{0.4pt}\\
(к защите / снять с защиты)\\[14pt]

\raggedleft
Научный руководитель:\\[10pt]
\rule{8cm}{0.4pt}\\[4pt]
\rule{8cm}{0.4pt}\\[14pt]

\raggedright
Проверка на объем заимствований:\\
\rule{4cm}{0.4pt}~\% авторского текста\\[10pt]

\raggedleft
<<Оценка>> \rule{2.5cm}{0.4pt}/\rule{2.5cm}{0.4pt}\\
(итоговая сумма баллов / оценка)\\

\vfill
\centering
Москва --- 2026 год
\end{titlepage}
\setcounter{page}{2}

%=====================================================================
%                          ОГЛАВЛЕНИЕ
%=====================================================================
\renewcommand{\contentsname}{ОГЛАВЛЕНИЕ}
\tableofcontents
\newpage

%=====================================================================
%                           ВВЕДЕНИЕ
%=====================================================================
\chapter*{ВВЕДЕНИЕ}
\addcontentsline{toc}{chapter}{ВВЕДЕНИЕ}

Вопрос о том, какие точки кривой имеют рациональные координаты, на первый
взгляд принадлежит теории чисел: речь идёт о решении уравнений в целых или
рациональных числах. Однако для кривых второго порядка ответ на него даёт
чисто геометрическая конструкция~--- проведение секущих через одну известную
точку кривой. Эта конструкция переводит арифметическую задачу в задачу о
пересечении прямой и коники и тем самым полностью её решает.

Работа посвящена этому переходу. Главный её результат~---
утверждение о том, что неособая коника, имеющая хотя бы одну
рациональную точку, устроена как проективная прямая: между её рациональными
точками и рациональными числами (с добавлением бесконечно удалённой точки)
существует взаимно однозначное соответствие, задаваемое пучком прямых.
Классическая задача о пифагоровых тройках получается как частный случай, когда
коника~--- окружность.

\textbf{Актуальность темы} определяется тем, что рациональная параметризация
коник служит наиболее доступным входом в современную арифметическую геометрию.
Тот же метод секущих, применённый к кривым третьего порядка, приводит к
групповому закону на эллиптической кривой~--- объекту, вокруг которого строится
значительная часть исследований последних десятилетий. Кроме того, материал
темы допускает наглядное изложение и потому пригоден для работы со школьниками,
интересующимися математикой.

\textbf{Объект исследования}~--- кривые второго порядка, заданные уравнениями
с рациональными коэффициентами.

\textbf{Предмет исследования}~--- метод секущих как способ описания множества
рациональных точек такой кривой и его проективная интерпретация.

\textbf{Цель работы}~--- изложить рациональную параметризацию конических
сечений, показать её проективную природу и продемонстрировать границы
применимости метода.

Для достижения цели поставлены следующие задачи:

\begin{enumerate}[leftmargin=*,topsep=2pt,itemsep=1pt]
\item выделить и доказать лемму о втором корне квадратного уравнения, лежащую
      в основании всего метода;
\item разобрать метод секущих на примере окружности и получить описание всех
      примитивных пифагоровых троек;
\item доказать основную теорему о том, что одна рациональная точка на неособой
      конике влечёт существование бесконечного их числа, и что секущие дают их все;
\item показать на примере, что коника с рациональными коэффициентами может не
      иметь ни одной рациональной точки, и указать критерий, отвечающий на вопрос
      о её существовании;
\item перевести конструкцию на язык проективной геометрии и установить её связь
      с теоремой Штейнера;
\item описать, что происходит с методом при переходе к кривым третьего порядка,
      и указать, почему для них аналогичное полное описание невозможно;
\item разобрать решения типовых задач по теме работы.
\end{enumerate}

\textbf{Методы исследования:} анализ и систематизация математической литературы,
метод секущих, аппарат однородных координат и проективных преобразований,
элементарные методы теории чисел (сравнения по модулю, бесконечный спуск).

\textbf{О доказательствах.} Утверждения, допускающие элементарное обоснование
доступными в курсе геометрии средствами, доказаны в работе полностью и
самостоятельно; к ним относятся лемма~\ref{lem:root}, теоремы~\ref{thm:circle},
\ref{thm:triples}, \ref{thm:main}, \ref{thm:three}, \ref{thm:proj},
\ref{thm:equiv} и следствие~\ref{cor:dich}. Результаты, строгое доказательство
которых требует аппарата, выходящего за рамки настоящей работы (теорема
Хассе~--- Минковского, теорема Штейнера, ассоциативность группового закона,
теорема Морделла), приведены со ссылками на источники и используются
как готовые факты; в тексте это каждый раз оговаривается.

\textbf{Структура работы.} Работа состоит из введения, пяти глав, заключения и
списка литературы. Первая глава носит вводный характер и разбирает основную
конструкцию на одном примере. Вторая переносит её на произвольную конику и
обсуждает вопрос существования рациональной точки. Третья даёт проективное
объяснение того, почему конструкция работает. Четвёртая показывает, где метод
перестаёт быть исчерпывающим. Пятая содержит разбор решений задач по материалу
предыдущих глав.

%=====================================================================
\chapter{РАЦИОНАЛЬНЫЕ ТОЧКИ НА ОКРУЖНОСТИ}
%=====================================================================

\section{Пифагоровы тройки: постановка задачи}

\begin{defn}
Пифагоровой тройкой называется тройка натуральных чисел $(a,b,c)$,
удовлетворяющая равенству $a^2+b^2=c^2$. Тройка называется примитивной, если
наибольший общий делитель её членов равен единице.
\end{defn}

Простейшие примеры известны с древности: $(3,4,5)$, $(5,12,13)$, $(8,15,17)$.
Задача состоит в том, чтобы описать все такие тройки. Если тройка $(a,b,c)$
не примитивна и $d$~--- общий делитель, то $(a/d,\,b/d,\,c/d)$ также является
пифагоровой тройкой. Поэтому достаточно описать примитивные тройки, а все
остальные получаются из них умножением на натуральное число.

Переведём задачу на геометрический язык. Разделим равенство $a^2+b^2=c^2$ на $c^2$. Получим
\[
  \left(\frac{a}{c}\right)^{\!2}+\left(\frac{b}{c}\right)^{\!2}=1 .
\]
Это означает, что точка с координатами $(a/c,\;b/c)$ лежит на единичной
окружности и обе её координаты рациональны. Обратно, если точка $(x,y)$ с
положительными рациональными координатами лежит на окружности $x^2+y^2=1$, то,
приведя $x$ и $y$ к общему знаменателю $c$, получим $x=a/c$, $y=b/c$ с
натуральными $a$, $b$ и, следовательно, $a^2+b^2=c^2$. Остальные рациональные
точки окружности получаются из точек первой четверти сменой знаков координат,
а также дают вырожденные решения с нулевым членом.

Значит, задача о пифагоровых тройках равносильна задаче об описании
рациональных точек окружности $x^2+y^2=1$; в этой форме мы её и решим.

\section{Лемма о втором корне и метод секущих}

Сначала докажем простой алгебраический факт, на котором основана вся
дальнейшая конструкция. Специальные свойства окружности при этом не нужны.

\begin{lem}[о втором корне]\label{lem:root}
Пусть квадратное уравнение $Ax^2+Bx+D=0$ с рациональными коэффициентами,
$A\neq 0$, имеет рациональный корень $x_0$. Тогда и второй его корень
рационален и равен $x_1=-B/A-x_0$.
\end{lem}

\begin{proof}
По теореме Виета сумма корней уравнения равна $-B/A$. Так как один из корней
равен $x_0$, второй равен $x_1=-B/A-x_0$. Величина $-B/A$ рациональна как
частное рациональных чисел при $A\neq 0$, число $x_0$ рационально по условию,
а разность двух рациональных чисел рациональна. Следовательно, $x_1$
рационально.
\end{proof}

Смысл леммы для дальнейшего таков: если прямая с рациональным угловым
коэффициентом проходит через рациональную точку кривой второго порядка, то
вторая точка их пересечения автоматически оказывается рациональной. Одна
известная точка порождает бесконечно много новых. Проведём это рассуждение
подробно для окружности.

Окружность $x^2+y^2=1$ содержит очевидные рациональные точки: $(\pm1,0)$ и
$(0,\pm1)$. Возьмём одну из них, например $A(-1,0)$, и будем проводить через
неё всевозможные прямые.

Прямая, проходящая через $A$ и не являющаяся вертикальной, имеет уравнение
$y=t(x+1)$, где $t$~--- её угловой коэффициент. Подставляя $y=t(x+1)$ в
уравнение окружности, получаем
\[
  x^2+t^2(x+1)^2=1,\qquad\text{или}\qquad (x-1)(x+1)+t^2(x+1)^2=0 .
\]
Вынося общий множитель, имеем
\[
  (x+1)\bigl[(x-1)+t^2(x+1)\bigr]=0 .
\]
Корень $x=-1$ отвечает самой точке $A$. Второй корень находится из уравнения
$x(1+t^2)=1-t^2$, откуда
\begin{equation}\label{eq:circle}
  x=\frac{1-t^2}{1+t^2},\qquad y=t(x+1)=\frac{2t}{1+t^2}.
\end{equation}

Построение показано на рис.~\ref{fig:secant}. Каждому рациональному значению
параметра $t$ отвечает своя секущая, а ей~--- своя рациональная точка
окружности; отмеченные на чертеже точки $M_1$ и $M_2$ дают соответственно
тройки $(3,4,5)$ и $(5,12,13)$.

\begin{figure}[htbp]
\centering
\begin{tikzpicture}[scale=2.6,>=Stealth,
  ln/.style={gray,line width=0.5pt},
  pt/.style={circle,fill=black,inner sep=1.3pt}]

  % оси
  \draw[->,line width=0.5pt] (-1.65,0) -- (2.5,0) node[below left] {$x$};
  \draw[->,line width=0.5pt] (0,-1.35) -- (0,1.55) node[below left] {$y$};
  \draw (0.03,1) -- (-0.03,1) node[left,scale=0.8] {$1$};

  % окружность
  \draw[line width=0.9pt] (0,0) circle (1);

  % касательная в A
  \draw[gray,dashed,line width=0.5pt] (-1,-1.2) -- (-1,1.3);
  \node[gray,scale=0.72,align=right,anchor=north east] at (-1.05,-0.95)
        {касательная в $A$\\ ($t=\infty$)};

  % секущие: t = 1/2, 2/3, 1
  \draw[ln] (-1,0) -- (0.82,0.91);            % t=1/2  -> (3/5,4/5)
  \draw[ln] (-1,0) -- (0.60,1.07);            % t=2/3  -> (5/13,12/13)
  \draw[ln] (-1,0) -- (0.22,1.22);            % t=1    -> (0,1)

  % точки
  \node[pt] (A) at (-1,0) {};
  \node[anchor=south east,scale=0.85] at (-1.02,0.05) {$A(-1;0)$};
  \node[pt] (M1) at (0.6,0.8) {};
  \node[pt] (M2) at (0.3846,0.9231) {};
  \node[pt] (M3) at (0,1) {};

  % выноски
  \draw[gray!55,line width=0.35pt] (M3) -- (0.50,1.38);
  \node[anchor=west,scale=0.8] at (0.50,1.38) {$M_3(0;\,1),\; t=1$};
  \draw[gray!55,line width=0.35pt] (M2) -- (1.02,1.05);
  \node[anchor=west,scale=0.8] at (1.02,1.05) {$M_2\!\left(\tfrac{5}{13};\tfrac{12}{13}\right)\!,\; t=\tfrac23$};
  \draw[gray!55,line width=0.35pt] (M1) -- (1.15,0.50);
  \node[anchor=west,scale=0.8] at (1.15,0.50) {$M_1\!\left(\tfrac35;\tfrac45\right)\!,\; t=\tfrac12$};
\end{tikzpicture}
\caption{Метод секущих на окружности $x^2+y^2=1$}
\label{fig:secant}
\end{figure}

\begin{thm}\label{thm:circle}
Формулы \eqref{eq:circle} задают взаимно однозначное соответствие между
множеством рациональных чисел $t$ и множеством рациональных точек окружности
$x^2+y^2=1$, отличных от точки $A(-1,0)$.
\end{thm}

\begin{proof}
Пусть сначала $t$ рационально. Знаменатель $1+t^2$ положителен и рационален,
поэтому правые части обеих формул рациональны, и соответствующая точка
окружности рациональна. Принадлежность окружности проверяется подстановкой:
\[
  \frac{(1-t^2)^2+(2t)^2}{(1+t^2)^2}
  =\frac{1-2t^2+t^4+4t^2}{(1+t^2)^2}
  =\frac{(1+t^2)^2}{(1+t^2)^2}=1 .
\]

Обратно, пусть $P(x_0,y_0)$~--- рациональная точка окружности, отличная от $A$.
Прямая $AP$ не вертикальна: единственная точка окружности с абсциссой $-1$ есть
сама $A$. Значит, её угловой коэффициент $t=y_0/(x_0+1)$ определён и рационален
как частное рациональных чисел. Прямая $y=t(x+1)$ пересекает окружность в
точках с абсциссами $-1$ и $(1-t^2)/(1+t^2)$, как показано выше, а точка $P$
лежит на этой прямой и отлична от $A$. Следовательно, $P$ есть точка
\eqref{eq:circle}, отвечающая данному рациональному $t$.

Остаётся проверить взаимную однозначность. Различным значениям $t$ отвечают
различные прямые пучка с центром $A$, а каждая из них пересекает окружность
ровно в одной точке, отличной от $A$. Следовательно, различным $t$ отвечают
различные точки, и соответствие взаимно однозначно.
\end{proof}

Заметим, что в доказательстве не использовались какие-либо особые свойства
окружности: нужны были только
лемма~\ref{lem:root} и тот факт, что прямая пересекает кривую второго порядка
не более чем в двух точках. Это наблюдение будет обобщено в главе~2 и получит
окончательную форму в главе~3.

Вертикальная прямая $x=-1$ в параметризацию не попадает: она касается
окружности в точке $A$. Удобно считать, что ей отвечает <<бесконечное>>
значение параметра; тогда соответствие становится взаимно однозначным между
окружностью и прямой, дополненной одной точкой. Этот приём будет формализован
в \S~3.1.

\section{Описание всех примитивных троек}

Запишем рациональный параметр в виде несократимой дроби $t=q/p$, где $p$ и $q$~---
целые, $p>0$, $\gcdru(p,q)=1$. Подстановка в формулы
\eqref{eq:circle} даёт
\[
  x=\frac{p^2-q^2}{p^2+q^2},\qquad y=\frac{2pq}{p^2+q^2},
\]
откуда получается тройка
\[
  a=p^2-q^2,\qquad b=2pq,\qquad c=p^2+q^2 .
\]
Непосредственная проверка подтверждает равенство $a^2+b^2=c^2$:
\[
  (p^2-q^2)^2+(2pq)^2=p^4-2p^2q^2+q^4+4p^2q^2=(p^2+q^2)^2 .
\]

\begin{thm}\label{thm:triples}
Пусть $p>q>0$, $\gcdru(p,q)=1$ и числа $p$ и $q$ имеют разную
чётность. Тогда тройка $(p^2-q^2,\;2pq,\;p^2+q^2)$ примитивна. Обратно, всякая
примитивная пифагорова тройка с чётным вторым членом представима в этом виде,
причём пара $(p,q)$ определена однозначно.
\end{thm}

\begin{proof}
Докажем сначала примитивность. Пусть $d$~--- общий делитель чисел
$a=p^2-q^2$ и $c=p^2+q^2$. Тогда $d$ делит их сумму $2p^2$ и их разность
$2q^2$. Поскольку $\gcdru(p,q)=1$, числа $p^2$ и $q^2$ взаимно
просты, и потому $d$ делит $2$. Так как $p$ и $q$ имеют разную чётность, число
$c=p^2+q^2$ нечётно, откуда $d$ нечётно, а значит $d=1$. Итак,
$\gcdru(a,c)=1$, и тем более наибольший общий делитель всех трёх
членов равен единице.

Докажем обратное утверждение. Пусть $(a,b,c)$~--- примитивная тройка. Числа $a$
и $b$ не могут быть оба чётными, иначе тройка не примитивна. Они не могут быть
и оба нечётными: в этом случае $a^2+b^2\equiv 2\pmod 4$, тогда как квадрат
целого числа сравним с $0$ или $1$ по модулю~4, и сумма двух квадратов не может
быть сравнима с~2 по модулю~4 при нечётных слагаемых. Следовательно, ровно одно
из чисел $a$, $b$ чётно; пусть по условию это $b$.

Тогда точка $(a/c,\;b/c)$~--- рациональная точка окружности, отличная от
$A(-1,0)$, и по теореме~\ref{thm:circle} ей отвечает единственное рациональное
значение $t=\dfrac{b/c}{a/c+1}=\dfrac{b}{a+c}$. Поскольку $a,b,c>0$, имеем
$0<t<1$: действительно, $b<a+c$, так как $b<c$. Запишем $t$ в виде несократимой
дроби $q/p$; тогда $p>q>0$ и, как показано выше,
\[
  \frac{a}{c}=\frac{p^2-q^2}{p^2+q^2},\qquad \frac{b}{c}=\frac{2pq}{p^2+q^2}.
\]

Покажем, что $p$ и $q$ имеют разную чётность. Оба чётными они быть не могут
ввиду $\gcdru(p,q)=1$. Пусть оба нечётны. Тогда $p^2+q^2\equiv2\pmod4$, а числа
$p^2-q^2$ и $2pq$ чётны, причём $pq$ нечётно; значит,
\[
  \frac{b}{c}=\frac{pq}{(p^2+q^2)/2},
\]
и эта дробь несократима, поскольку всякий общий простой делитель чисел $pq$ и
$(p^2+q^2)/2$ делил бы $p$ или $q$, а вместе с $p^2+q^2$~--- и второе из них. Дробь
$b/c$ также несократима: $\gcdru(b,c)=1$ ввиду примитивности тройки (общий
делитель $b$ и $c$ делил бы и $a^2$). По единственности несократимой записи
$b=pq$ нечётно, что противоречит чётности $b$.

Итак, $p$ и $q$ разной чётности, и по доказанному в первой части дробь
$(p^2-q^2)/(p^2+q^2)$ несократима. Несократима и дробь $a/c$, так как
$\gcdru(a,c)=1$. По единственности несократимой записи рационального числа
$a=p^2-q^2$, $c=p^2+q^2$, а тогда и $b=2pq$. Однозначность пары $(p,q)$ следует
из того, что $t=q/p$ определено точкой $(a/c,b/c)$ однозначно
(теорема~\ref{thm:circle}), а несократимая запись $t$ единственна.
\end{proof}

Другие подходы к описанию пифагоровых троек см. в \cite{hatcher}. Приведём
первые тройки, получаемые по этим формулам:
{\small\[
\begin{array}{lll}
p=2,\;q=1:\;(3,4,5); & p=3,\;q=2:\;(5,12,13); & p=4,\;q=1:\;(15,8,17);\\[2pt]
p=4,\;q=3:\;(7,24,25); & p=5,\;q=2:\;(21,20,29). &
\end{array}
\]}
Условие разной чётности существенно. При $p=3$, $q=1$ формулы дают $(8,6,10)$~---
тройку, не являющуюся примитивной; она получается удвоением тройки $(4,3,5)$.

\section{Секущие как стереографическая проекция}

Конструкцию из \S~1.2 можно истолковать геометрически. Параметр $t$ есть
угловой коэффициент прямой, проходящей через $A(-1,0)$. Эта же прямая
пересекает ось ординат в точке $(0,t)$. Тем самым каждой точке окружности,
отличной от $A$, сопоставляется точка оси ординат, и это сопоставление взаимно
однозначно.

Такое отображение называется стереографической проекцией окружности на прямую
из точки $A$. Оно переводит окружность без одной точки в прямую; сама точка
$A$, как уже отмечалось, отвечает <<бесконечно удалённой>> точке прямой.
Важно, что проекция задана формулами с рациональными коэффициентами в обе
стороны, а потому рациональные точки окружности переходят ровно в рациональные
точки прямой.

Это свойство (рациональность отображения и обратного к нему) лежит в основе
определения рациональной эквивалентности в \S~3.5. В терминах этого параграфа
окружность с выколотой точкой рационально эквивалентна прямой, поэтому
множества их рациональных точек устроены одинаково.

%=====================================================================
\chapter{КОНИКИ НАД ПОЛЕМ РАЦИОНАЛЬНЫХ ЧИСЕЛ}
%=====================================================================

\section{Кривые второго порядка и рациональные точки}

\begin{defn}
Коникой (кривой второго порядка) над полем рациональных чисел называется
множество точек плоскости, координаты которых удовлетворяют уравнению
\begin{equation}\label{eq:conic}
  ax^2+bxy+cy^2+dx+ey+f=0,
\end{equation}
где коэффициенты $a,b,c,d,e,f$ рациональны и хотя бы один из коэффициентов
$a,b,c$ отличен от нуля.
\end{defn}

Кривой сопоставляется симметрическая матрица
\[
  M=\begin{pmatrix}
     a & b/2 & d/2\\
     b/2 & c & e/2\\
     d/2 & e/2 & f
    \end{pmatrix}.
\]

\begin{defn}\label{def:nonsing}
Коника называется неособой (невырожденной), если определитель матрицы $M$
отличен от нуля \cite[с.~118]{atanasyan}.
\end{defn}

Особые коники распадаются на пару прямых (возможно, совпавших или мнимых) либо
вырождаются в точку. Их рациональные точки описываются тривиально, поэтому
далее рассматриваются только неособые коники.

\begin{defn}
Рациональной точкой коники называется точка кривой, обе координаты которой
рациональны. Множество рациональных точек коники $C$ обозначается $C(\Q)$.
\end{defn}

Возникают два вопроса. Первый: пусто ли множество $C(\Q)$? Второй: если оно
непусто, как его описать? Оказывается, что второй вопрос решается полностью и
элементарно, а первый требует существенно более тонких средств. Начнём со
второго.

\section{Основная теорема о рациональной параметризации}

Сначала убедимся, что прямая и коника не могут иметь слишком много общих точек.
Это утверждение используется на протяжении всей работы, и его удобно
зафиксировать отдельно.

\begin{lem}\label{lem:two}
Прямая, не лежащая целиком на конике, имеет с ней не более двух общих точек.
\end{lem}

\begin{proof}
Пусть прямая задана параметрически: $x=x_0+\alpha s$, $y=y_0+\beta s$, где
$(\alpha,\beta)\neq(0,0)$. Подставим эти выражения в уравнение коники. Каждое
слагаемое левой части есть многочлен от $s$ степени не выше второй, поэтому и
вся левая часть есть многочлен от $s$ степени не выше второй. Если он
тождественно равен нулю, прямая целиком лежит на конике, что исключено условием.
В противном случае это ненулевой многочлен степени не выше двух, и он имеет не
более двух корней. Каждому общему значению параметра отвечает ровно одна точка
прямой, поэтому общих точек не более двух.
\end{proof}

\begin{thm}\label{thm:main}
Пусть $C$~--- неособая коника с рациональными коэффициентами и
$P_0\in C(\Q)$. Тогда множество $C(\Q)$ бесконечно, и всякая его точка,
отличная от $P_0$, получается как вторая точка пересечения коники с прямой,
проходящей через $P_0$ и имеющей рациональный угловой коэффициент (либо с
вертикальной прямой, проходящей через $P_0$).
\end{thm}

\begin{proof}
Пусть $P_0=(x_0,y_0)$ и пусть $\ell$~--- прямая через $P_0$ с угловым
коэффициентом $t$, то есть $y=y_0+t(x-x_0)$. Подставим это выражение в
уравнение коники. Все слагаемые при этом становятся многочленами от $x$ степени
не выше второй, и уравнение принимает вид
\[
  A(t)\,x^2+B(t)\,x+D(t)=0,
\]
где $A$, $B$, $D$~--- многочлены от $t$ с рациональными коэффициентами; их
рациональность следует из рациональности коэффициентов коники и параметра $t$.
Число $x_0$ является корнем этого уравнения по построению, поскольку точка
$P_0$ лежит и на прямой, и на конике.

Пусть $A(t)\neq 0$. Тогда по лемме~\ref{lem:root} второй корень рационален и
равен
\[
  x_1=-\frac{B(t)}{A(t)}-x_0,
\]
а соответствующая ордината $y_1=y_0+t(x_1-x_0)$ рациональна как результат
рациональных операций над рациональными числами. Итак, каждому рациональному
$t$, для которого $A(t)\neq0$, отвечает рациональная точка коники.

Покажем, что различным значениям параметра отвечают, за конечным числом
исключений, различные точки. Прямые с различными угловыми коэффициентами,
проходящие через одну точку, различны и пересекаются только в этой точке $P_0$.
Поэтому их вторые точки пересечения с коникой либо различны, либо совпадают с
$P_0$. Последнее возможно лишь для касательной к конике в точке $P_0$, то есть
не более чем для одного значения $t$. Поскольку рациональных чисел бесконечно
много, а исключительных значений конечное число (не более двух корней
многочлена $A$ и одно значение для касательной), получается бесконечно много
различных рациональных точек. Значит, множество $C(\Q)$ бесконечно.

Обратно, пусть $Q\in C(\Q)$, $Q\neq P_0$. Прямая $P_0Q$ проходит через две точки
с рациональными координатами. Если их абсциссы различны, её угловой коэффициент
$t=(y_Q-y_0)/(x_Q-x_0)$ рационален; если абсциссы совпадают, прямая вертикальна.
По лемме~\ref{lem:two} прямая $P_0Q$ пересекает конику не более чем в двух
точках, и обе они здесь известны: $P_0$ и $Q$. Значит, $Q$ есть в точности
вторая точка пересечения коники с прямой указанного вида, то есть получается
описанной конструкцией.
\end{proof}

Разберём исключительные случаи, упомянутые в доказательстве.

Во-первых, при некоторых значениях $t$ может оказаться $A(t)=0$. Геометрически
это означает, что прямая $\ell$ параллельна асимптотическому направлению коники
и пересекает её лишь в одной конечной точке; вторая точка пересечения уходит на
бесконечность. Таких значений $t$ не более двух, поскольку $A$~--- многочлен
степени не выше второй.

Во-вторых, одному значению $t$ отвечает касательная к конике в точке $P_0$: для
неё обе точки пересечения совпадают с $P_0$, и новой точки не получается.

Оба исключения исчезают при переходе к проективной плоскости, где бесконечно
удалённые точки равноправны с конечными. Это ещё один довод в пользу
проективной формулировки, которой посвящена глава~3.

\begin{cor}\label{cor:dich}
Для неособой коники с рациональными коэффициентами возможны ровно две ситуации:
либо $C(\Q)$ пусто, либо $C(\Q)$ бесконечно. Промежуточного случая~--- конечного
ненулевого числа рациональных точек~--- не существует.
\end{cor}

\begin{proof}
Если $C(\Q)$ непусто, в нём найдётся точка $P_0$, и по теореме~\ref{thm:main}
множество $C(\Q)$ бесконечно. Значит, конечное ненулевое число рациональных
точек невозможно.
\end{proof}

Это утверждение резко отличает коники от кривых более высокого порядка, для
которых конечное ненулевое число рациональных точек~--- обычное явление.
Соответствующий пример разобран в упражнении~\ref{ex:cubic}, а причина различия
обсуждается в главе~4.

\section{Коники без рациональных точек}

Следствие~\ref{cor:dich} оставляет открытым вопрос: как узнать, какая из двух
ситуаций имеет место? Приведём пример, показывающий, что первая из них
действительно реализуется.

\begin{thm}\label{thm:three}
Окружность $x^2+y^2=3$ не содержит ни одной рациональной точки.
\end{thm}

\begin{proof}
Предположим противное: пусть $(x,y)$~--- рациональная точка. Приведя координаты
к общему знаменателю, запишем $x=X/Z$, $y=Y/Z$, где $X$, $Y$, $Z$~--- целые,
$Z\neq0$, и $\gcdru(X,Y,Z)=1$. Подстановка и умножение на $Z^2$ дают
\[
  X^2+Y^2=3Z^2 .
\]
Рассмотрим это равенство по модулю~3. Квадрат целого числа сравним с $0$ или
$1$ по модулю~3: действительно, при $n\equiv0$ имеем $n^2\equiv0$, при
$n\equiv\pm1$ имеем $n^2\equiv1$. Правая часть равенства сравнима с нулём,
поэтому $X^2+Y^2\equiv0\pmod 3$. Перебор всех возможностей для пары остатков
($0+0$, $0+1$, $1+0$, $1+1$) показывает, что сумма сравнима с нулём только в
случае $0+0$. Следовательно, $X^2\equiv Y^2\equiv 0 \pmod 3$, а поскольку 3
просто, отсюда $X\equiv Y\equiv 0\pmod 3$.

Запишем $X=3X_1$, $Y=3Y_1$. Подстановка даёт $9X_1^2+9Y_1^2=3Z^2$, откуда
$3(X_1^2+Y_1^2)=Z^2$. Значит, 3 делит $Z^2$, и, поскольку число 3 простое, 3 делит $Z$.
Но тогда числа $X$, $Y$, $Z$ имеют общий делитель~3, что противоречит
предположению о взаимной простоте.
\end{proof}

Изложенное рассуждение (о сравнениях и квадратичных вычетах
см.~\cite{vinogradov}) является простейшим вариантом метода бесконечного
спуска: из существования решения выводится существование решения
<<меньшего>>, и этот процесс не может продолжаться бесконечно. Аналогичное
рассуждение для окружности $x^2+y^2=7$ разобрано в упражнении~\ref{ex:seven}.

Общий ответ на вопрос о существовании рациональной точки даёт следующий
результат. Его доказательство опирается на теорию $p$-адических чисел и выходит
за рамки настоящей работы, поэтому теорема приводится со ссылкой и используется
далее как готовый факт.

\begin{thm}[принцип Хассе~--- Минковского для коник]\label{thm:hasse}
Неособая коника с рациональными коэффициентами имеет рациональную точку тогда и
только тогда, когда она имеет точку над полем вещественных чисел и над полем
$p$-адических чисел для каждого простого $p$ \cite[гл.~IV, \S~3]{serre}.
Доказательство см. также в \cite[гл.~1, \S~7]{borevich}.
\end{thm}

Смысл теоремы в том, что глобальная разрешимость равносильна совокупности локальных разрешимостей: препятствий, не
обнаруживаемых ни вещественными соображениями, ни сравнениями по модулям
степеней простых, для коник не существует. Практически проверка сводится к
конечному числу условий, поскольку для всех нечётных простых, не делящих
определитель матрицы коники (после приведения коэффициентов к целым), локальная
разрешимость выполняется автоматически; отдельно проверяются лишь простое
$p=2$, простые делители определителя и вещественный случай. Для кривых третьего порядка
аналогичное утверждение, вообще говоря, неверно; известны контрпримеры,
простейший из которых принадлежит Сельмеру \cite[гл.~IV]{silverman}.

\section{Примеры: уравнение Пелля и треугольники с углом 120\textdegree}

\begin{exmp}
Рассмотрим гиперболу $x^2-2y^2=1$. Она содержит очевидную рациональную точку
$P_0(1,0)$. Проведём через неё прямую $y=t(x-1)$. Подстановка даёт
\[
  x^2-2t^2(x-1)^2=1,\qquad\text{то есть}\qquad (x-1)(x+1)=2t^2(x-1)^2 .
\]
Сокращая на $x-1$ (что законно при $x\neq1$), получаем $x+1=2t^2(x-1)$, откуда
\[
  x=\frac{2t^2+1}{2t^2-1},\qquad y=\frac{2t}{2t^2-1}.
\]
При $t=1$ формулы дают точку $(3,2)$, при $t=3/4$~--- точку $(17,12)$, при
$t=5/7$~--- точку $(99,70)$. Это решения уравнения Пелля
$x^2-2y^2=1$ в целых числах.
\end{exmp}

Отметим, однако, что метод секущих описывает все рациональные точки
гиперболы, но вопрос о том, какие из них целочисленны, он не решает. Целые
решения уравнения Пелля образуют последовательность, задаваемую степенями числа
$3+2\sqrt2$, и их описание требует иных средств~--- теории непрерывных дробей
или арифметики квадратичных полей \cite[гл.~2]{borevich}. Рациональная параметризация есть решение задачи о
рациональных точках, и не более того.

\begin{exmp}
Рассмотрим эллипс $x^2+xy+y^2=1$. Он связан с задачей о треугольниках, у которых
один угол равен $120^\circ$, а стороны целочисленны. Действительно, по теореме
косинусов для треугольника со сторонами $a$, $b$ и углом $120^\circ$ между ними
противолежащая сторона $c$ удовлетворяет равенству
\[
  c^2=a^2+b^2-2ab\cos 120^\circ=a^2+ab+b^2 .
\]
Разделив на $c^2$, получаем, что точка $(a/c,\;b/c)$ лежит на указанном
эллипсе. Проведём через рациональную точку $(1,0)$ прямую $y=t(x-1)$.
Подстановка в уравнение эллипса даёт
\[
  x^2+tx(x-1)+t^2(x-1)^2=1,\qquad\text{то есть}\qquad
  (x-1)\bigl[(x+1)+tx+t^2(x-1)\bigr]=0,
\]
откуда $x(1+t+t^2)=t^2-1$ и
\[
  x=\frac{t^2-1}{t^2+t+1},\qquad y=-\,\frac{t(t+2)}{t^2+t+1}.
\]
\end{exmp}

Положим $t=n/m$ с натуральными $m>n$. Подстановка даёт точку
\[
  \left(-\frac{m^2-n^2}{m^2+mn+n^2},\;-\frac{2mn+n^2}{m^2+mn+n^2}\right),
\]
а поскольку эллипс симметричен относительно начала координат, вместе с ней на
нём лежит и точка с противоположными координатами. Отсюда получаются
классические формулы, которые проверяются и непосредственной подстановкой: при
взаимно простых натуральных $m>n$
\[
  a=m^2-n^2,\qquad b=2mn+n^2,\qquad c=m^2+mn+n^2 .
\]
Проверка:
$a^2+ab+b^2=m^4+2m^3n+3m^2n^2+2mn^3+n^4=(m^2+mn+n^2)^2=c^2$.

При $m=2$, $n=1$ получается треугольник со сторонами $(3,5,7)$, при $m=3$,
$n=1$~--- треугольник $(8,7,13)$, при $m=4$, $n=1$~--- треугольник $(15,9,21)$,
кратный тройке $(5,3,7)$. Такие тройки называются эйзенштейновыми. Вопрос о
примитивности решается рассуждением, аналогичным теореме~\ref{thm:triples}, но
проводимым по модулю~3: при взаимно простых $m>n$ тройка примитивна тогда и
только тогда, когда $m-n$ не делится на~3 (в примере $m=4$, $n=1$ это условие
нарушено).

Сопоставление двух примеров показывает главное: конструкция не зависит от того,
является ли коника эллипсом или гиперболой (для параболы рассуждение то же). Все различия между
ними~--- аффинные, тогда как метод секущих по своей природе проективен. Это
соображение и мотивирует переход к следующей главе.

%=====================================================================
\chapter{ПРОЕКТИВНЫЙ ВЗГЛЯД}
%=====================================================================

\section{Проективная плоскость и однородные координаты}

\begin{defn}
Проективной плоскостью над полем $k$ называется множество классов
эквивалентности ненулевых троек $(x_0:x_1:x_2)$ элементов поля $k$, где тройки
$(x_0,x_1,x_2)$ и $(\lambda x_0,\lambda x_1,\lambda x_2)$ при $\lambda\neq0$
отождествляются. Обозначение: $\PP^2(k)$. Числа $(x_0:x_1:x_2)$ называются
однородными координатами точки \cite[гл.~III, \S~1]{atanasyan}.
\end{defn}

Точки, у которых $x_0\neq0$, находятся во взаимно однозначном соответствии с
точками обычной (аффинной) плоскости по правилу
\[
  (x_0:x_1:x_2)\;\longmapsto\;\left(\frac{x_1}{x_0},\;\frac{x_2}{x_0}\right).
\]
Точки с $x_0=0$ называются бесконечно удалёнными; они образуют одну прямую,
называемую бесконечно удалённой прямой. Содержательно каждая такая точка
отвечает направлению на аффинной плоскости: параллельные прямые пересекаются в
общей бесконечно удалённой точке.

Проективная прямая $\PP^1(k)$ определяется аналогично как множество классов пар
$(x_0:x_1)$. Для $k=\Q$ она состоит из всех рациональных чисел и одной
дополнительной точки, обозначаемой $\infty$. Именно эта конструкция придаёт
точный смысл оговорке о <<бесконечном значении параметра>> из \S~1.2.

Точка проективной плоскости называется рациональной, если её можно записать
однородными координатами, все из которых рациональны. Это определение корректно:
умножение всех координат на общий знаменатель не меняет точку, поэтому у
рациональной точки всегда есть и запись целыми взаимно простыми координатами.

\section{Коника в проективной плоскости}

Уравнение коники \eqref{eq:conic} превращается в однородное подстановкой
$x=x_1/x_0$, $y=x_2/x_0$ и умножением на $x_0^2$:
\[
  a\,x_1^2+b\,x_1x_2+c\,x_2^2+d\,x_0x_1+e\,x_0x_2+f\,x_0^2=0 .
\]
Левая часть есть квадратичная форма от трёх переменных; её обращение в нуль
корректно определено на классах эквивалентности, поскольку при замене
$(x_0,x_1,x_2)$ на $(\lambda x_0,\lambda x_1,\lambda x_2)$ форма умножается на
$\lambda^2$, а обращение в нуль при этом сохраняется. Коника в $\PP^2$ есть
множество нулей неособой квадратичной формы; неособость означает, что
определитель матрицы формы отличен от нуля (определение~\ref{def:nonsing}).

Проективная точка зрения устраняет аффинные различия. Эллипс, парабола и
гипербола отличаются лишь тем, как коника расположена относительно бесконечно
удалённой прямой: эллипс с ней не пересекается, парабола касается, гипербола
пересекает в двух точках (рис.~\ref{fig:infinity}). Соответствующие выкладки
разобраны в упражнении~\ref{ex:infty}. Над полем
вещественных чисел как проективные кривые все три неотличимы: существует
проективное преобразование, переводящее одну в другую
\cite[гл.~II, \S~12]{efimov}. Над полем $\Q$ это уже неверно: например,
окружности $x^2+y^2=1$ и $x^2+y^2=3$ не переводятся друг в друга никаким
проективным преобразованием с рациональными коэффициентами, поскольку такое
преобразование сохраняет рациональность точек (теорема~\ref{thm:three}).

\begin{figure}[htbp]
\centering
\begin{tikzpicture}[scale=0.92]
% ---- эллипс ----
\begin{scope}[xshift=0cm]
  \draw[->,line width=0.5pt] (-1.9,0) -- (1.9,0) node[below left,scale=0.75] {$x$};
  \draw[->,line width=0.5pt] (0,-1.6) -- (0,1.6) node[below left,scale=0.75] {$y$};
  \draw[line width=0.9pt] (0,0) ellipse (1.2 and 0.85);
  \node[scale=0.8] at (0,2.0) {Эллипс $x^2+2y^2=1$};
  \node[scale=0.75,align=center] at (0,-2.05) {вещественных бесконечно\\удалённых точек нет};
\end{scope}
% ---- гипербола ----
\begin{scope}[xshift=5.1cm]
  \draw[->,line width=0.5pt] (-1.9,0) -- (1.9,0) node[below left,scale=0.75] {$x$};
  \draw[->,line width=0.5pt] (0,-1.6) -- (0,1.6) node[below left,scale=0.75] {$y$};
  \draw[gray,dashed,line width=0.5pt] (-1.75,-1.24) -- (1.75,1.24);
  \draw[gray,dashed,line width=0.5pt] (-1.75,1.24) -- (1.75,-1.24);
  \draw[line width=0.9pt,domain=-1.15:1.15,smooth,variable=\u]
        plot ({cosh(\u)},{sinh(\u)/1.4142});
  \draw[line width=0.9pt,domain=-1.15:1.15,smooth,variable=\u]
        plot ({-cosh(\u)},{sinh(\u)/1.4142});
  \node[scale=0.8] at (0,2.0) {Гипербола $x^2-2y^2=1$};
  \node[scale=0.75,align=center] at (0,-2.05)
       {две точки $(0:\pm\sqrt2:1)$,\\обе иррациональные};
\end{scope}
% ---- парабола ----
\begin{scope}[xshift=10.2cm]
  \draw[->,line width=0.5pt] (-1.9,0) -- (1.9,0) node[below left,scale=0.75] {$x$};
  \draw[->,line width=0.5pt] (0,-1.6) -- (0,1.6) node[below left,scale=0.75] {$y$};
  \draw[gray,dashed,line width=0.5pt] (0,-1.45) -- (0,1.5);
  \draw[line width=0.9pt,domain=-1.22:1.22,smooth,variable=\x]
        plot ({\x},{\x*\x});
  \node[scale=0.8] at (0,2.0) {Парабола $y=x^2$};
  \node[scale=0.75,align=center] at (0,-2.05)
       {одна точка $(0:0:1)$\\кратности два~--- касание};
\end{scope}
\end{tikzpicture}
\caption{Положение коник относительно бесконечно удалённой прямой}
\label{fig:infinity}
\end{figure}

\section{Пучок прямых и точки коники}

\begin{defn}
Пучком прямых с центром $P_0$ называется множество всех прямых проективной
плоскости, проходящих через точку $P_0$.
\end{defn}

Пучок прямых сам является проективной прямой: прямые пучка параметризуются
точками $\PP^1$. Это позволяет переформулировать теорему~\ref{thm:main} в
окончательном виде.

\begin{thm}\label{thm:proj}
Пусть $C$~--- неособая коника в $\PP^2(\Q)$ и $P_0\in C(\Q)$. Тогда отображение,
сопоставляющее каждой прямой $\ell$ пучка с центром $P_0$ вторую точку её
пересечения с $C$ (а касательной к $C$ в точке $P_0$~--- саму точку $P_0$),
является взаимно однозначным соответствием между пучком прямых и коникой $C$.
Это соответствие переводит рациональные прямые пучка ровно в рациональные точки
коники.
\end{thm}

\begin{proof}
Прежде всего, ни одна прямая не лежит целиком на неособой конике: в противном
случае квадратичная форма делилась бы на линейную, что означало бы вырождение
коники в пару прямых и обращение определителя матрицы формы в нуль. Значит, по
лемме~\ref{lem:two} каждая прямая пучка пересекает $C$ не более чем в двух
точках, и одна из них~--- точка $P_0$.

Проверим существование и единственность второй точки. Пусть $F$~--- квадратичная
форма коники, $\Phi$~--- соответствующая ей симметричная билинейная форма, так
что $F(X)=\Phi(X,X)$. Выберем на прямой $\ell$ пучка точку $R\neq P_0$; тогда
точки $\ell$ суть $\lambda P_0+\mu R$, $(\lambda:\mu)\in\PP^1$. Поскольку
$F(P_0)=0$,
\[
  F(\lambda P_0+\mu R)=2\lambda\mu\,\Phi(P_0,R)+\mu^2F(R)
  =\mu\bigl(2\lambda\,\Phi(P_0,R)+\mu F(R)\bigr).
\]
Эта бинарная форма не равна тождественно нулю, так как $\ell\not\subset C$.
Корень $\mu=0$ отвечает точке $P_0$, второй корень
$(\lambda:\mu)=\bigl(F(R):-2\Phi(P_0,R)\bigr)$ определён однозначно и даёт
единственную вторую точку пересечения
\[
  Q=F(R)\,P_0-2\Phi(P_0,R)\,R .
\]
Она совпадает с $P_0$ тогда и только тогда, когда $\Phi(P_0,R)=0$, то есть
когда $\ell$~--- касательная к $C$ в точке $P_0$; в этом случае по условию ей
сопоставляется сама точка $P_0$. Заметим, что в проективной записи случай
<<второй точки на бесконечности>> ничем не выделен.

Обратное соответствие строится так. Для точки $Q\in C$, отличной от $P_0$,
существует единственная прямая, проходящая через $P_0$ и $Q$, и она принадлежит
пучку; для $Q=P_0$ берётся касательная к $C$ в точке $P_0$, которая также
единственна. Построенные отображения взаимно обратны, поэтому соответствие
взаимно однозначно.

Наконец, рациональность в обе стороны. Если прямая $\ell$ рациональна, на ней
можно взять рациональную точку $R$, и тогда координаты $Q$ получаются из
координат $P_0$ и $R$ рациональными операциями, то есть $Q$ рациональна
(проективный вариант леммы~\ref{lem:root}). Обратно, две различные
рациональные точки $P_0$ и $Q$ определяют прямую с рациональными
коэффициентами, а касательная в рациональной точке $P_0$ задаётся уравнением
$\Phi(P_0,X)=0$ с рациональными коэффициентами.
\end{proof}

В проективной формулировке исчезают обе оговорки, сопровождавшие
теорему~\ref{thm:main}. Значения параметра, при которых ведущий коэффициент
обращался в нуль, теперь отвечают прямым, вторая точка пересечения которых
бесконечно удалена~--- но в $\PP^2$ такая точка ничем не хуже прочих.
Касательная получает свою точку. Соответствие становится полным.

\section{Связь с теоремой Штейнера}

Построенное соответствие не случайно и является частным проявлением
классического результата проективной геометрии. Его доказательство опирается на
теорию проективных соответствий и в настоящей работе не приводится; теорема
используется как известный факт.

\begin{thm}[Штейнер]\label{thm:steiner}
Пусть $A$ и $B$~--- две различные точки неособой коники $C$. Тогда отображение
пучка прямых с центром $A$ в пучок прямых с центром $B$, при котором прямой
$AX$ ставится в соответствие прямая $BX$ ($X\in C$), является проективным
преобразованием пучков. Обратно, если задано проективное, но не перспективное
соответствие двух пучков, то множество точек пересечения соответственных прямых
есть неособая коника \cite[гл.~II, \S~14]{efimov}. Изложение см. также в
\cite[гл.~8]{coxeter} и \cite[\S~20]{hartshorne}.
\end{thm}

Теорема Штейнера может быть принята за определение коники, не использующее
координат: коника есть геометрическое место точек пересечения соответственных
прямых двух проективно соответственных пучков. Доказанная выше
теорема~\ref{thm:proj} тесно связана с первой половиной этого утверждения: в ней
коника отождествляется с одним пучком, а соответствие Штейнера получается
композицией двух таких отождествлений~--- с центрами $A$ и~$B$.

Отсюда видно, почему метод секущих устойчив к замене окружности на произвольную
конику: он опирается не на метрические свойства кривой, а на проективный факт о
числе точек пересечения прямой и коники, то есть на лемму~\ref{lem:two}.

\section{Коника с рациональной точкой как проективная прямая}

\begin{defn}
Две кривые называются рационально эквивалентными над $\Q$, если между ними
существует взаимно однозначное соответствие, задаваемое в обе стороны формулами
с рациональными коэффициентами (за исключением конечного числа точек).
\end{defn}

\begin{thm}\label{thm:equiv}
Неособая коника над $\Q$, имеющая хотя бы одну рациональную точку, рационально
эквивалентна проективной прямой $\PP^1$. Коника, не имеющая рациональных точек,
проективной прямой над $\Q$ не эквивалентна.
\end{thm}

\begin{proof}
Пусть $P_0$~--- рациональная точка коники $C$. По теореме~\ref{thm:proj}
соответствие между пучком прямых с центром $P_0$ и коникой $C$ взаимно
однозначно, а формулы, его задающие, получены в теореме~\ref{thm:main} и
являются дробно-рациональными с рациональными коэффициентами в обе стороны: в
одну сторону это выражение второй точки пересечения через угловой коэффициент,
в другую~--- выражение углового коэффициента через координаты точки. Пучок
прямых есть проективная прямая, что и даёт требуемую эквивалентность.

Второе утверждение следует из того, что проективная прямая $\PP^1(\Q)$ содержит
бесконечно много рациональных точек, а формулы соответствия не определены лишь
в конечном числе точек. Если бы коника без рациональных точек была рационально
эквивалентна прямой, то образ рациональной точки прямой, в которой формулы
определены, был бы рациональной точкой коники, поскольку формулы соответствия
имеют рациональные коэффициенты. Полученное противоречие завершает
доказательство.
\end{proof}

Подведём итог первых трёх глав: с точки зрения арифметики неособая коника с рациональной точкой не отличается от прямой. Все её
рациональные точки перечисляются одним параметром, пробегающим рациональные
числа. Тем самым задача о рациональных точках на конике решена полностью: если
на конике есть хотя бы одна рациональная точка, их бесконечно много и все они
описываются одним параметром, а если нет ни одной, то их нет вовсе
(следствие~\ref{cor:dich}).

%=====================================================================
\chapter{ЧТО МЕНЯЕТСЯ ДЛЯ КРИВЫХ ТРЕТЬЕГО ПОРЯДКА}
%=====================================================================

\section{Метод хорд и касательных}

Естественно спросить, сохраняется ли развитая конструкция для кривых более
высокого порядка. Ответ отрицателен, и разбор причины проясняет, за счёт чего
работал метод секущих.

Рассмотрим кривую третьего порядка, приведённую к виду Вейерштрасса
\[
  y^2=x^3+ax+b,
\]
где $a$, $b$ рациональны, а кривая неособа, то есть дискриминант
$\Delta=-16(4a^3+27b^2)$ отличен от нуля \cite[гл.~I, \S~3]{silverman}.

\begin{prop}\label{prop:third}
Пусть прямая с рациональными коэффициентами пересекает кубику в трёх точках
(с учётом кратности) и две из них рациональны. Тогда рациональна и третья.
\end{prop}

\begin{proof}
Подставив уравнение прямой $y=kx+m$ с рациональными $k$ и $m$ в уравнение
кривой, получим кубическое уравнение относительно $x$:
\[
  x^3-k^2x^2+(a-2km)x+(b-m^2)=0,
\]
все коэффициенты которого рациональны. По формулам Виета сумма его корней равна
$k^2$, то есть рациональна. Два корня рациональны по условию, поэтому третий
корень равен разности рационального числа $k^2$ и суммы двух рациональных
чисел, а значит рационален. Соответствующая ордината $y=kx+m$ также рациональна.
\end{proof}

В этом главное отличие от коники. Там одна известная точка порождала все
остальные: прямая пересекала кривую ещё ровно в одной точке, и параметр пробегал
целую прямую. Здесь же для построения новой точки нужны две уже известные (или
одна с касательной в ней), а результат зависит от выбора пары. Такой процесс даёт
не параметризацию, а операцию над точками.

\section{Сложение точек}

Указанную операцию принято оформлять так, чтобы она превращала множество
рациональных точек в группу.

\begin{defn}\label{def:add}
Пусть $P$ и $Q$~--- точки кубики. Проведём через них прямую, найдём третью точку
пересечения $R$ и отразим её относительно оси абсцисс. Полученная точка
называется суммой $P+Q$. Если $P=Q$, вместо секущей берётся касательная в точке
$P$. Нейтральным элементом служит бесконечно удалённая точка $O$ кривой
\cite[гл.~I, \S~2]{silverman}.
\end{defn}

Построение суммы двух точек изображено на рис.~\ref{fig:cubic} слева: секущая
$PQ$ пересекает кривую в третьей точке $R$, и отражение этой точки относительно
оси абсцисс даёт $P+Q$. Справа на том же рисунке показан пример конечной группы,
разобранный в упражнении~\ref{ex:cubic}.

\begin{figure}[htbp]
\centering
\begin{tikzpicture}[scale=0.95,>=Stealth,pt/.style={circle,fill=black,inner sep=1.2pt}]
% ---- слева: сложение P+Q на y^2 = x^3 - 2x + 2 ----
\begin{scope}
  \draw[->,line width=0.5pt] (-2.1,0) -- (2.6,0) node[below left,scale=0.75] {$x$};
  \draw[->,line width=0.5pt] (0,-2.5) -- (0,2.7) node[below left,scale=0.75] {$y$};
  \draw (-1,0.06) -- (-1,-0.06) node[below,scale=0.7] {$-1$};
  \draw (1,0.06) -- (1,-0.06) node[below,scale=0.7] {$1$};
  \draw[line width=0.9pt,domain=-1.62:1.75,smooth,variable=\x]
        plot ({\x},{sqrt(max(\x*\x*\x-2*\x+2,0))});
  \draw[line width=0.9pt,domain=-1.62:1.75,smooth,variable=\x]
        plot ({\x},{-sqrt(max(\x*\x*\x-2*\x+2,0))});
  % секущая y = -0.36603 x + 1.36603
  \draw[gray,line width=0.6pt] (-1.55,1.933) -- (1.95,0.652);
  \node[gray,scale=0.72,anchor=north east] at (1.95,0.60) {секущая $PQ$};
  % отражение
  \draw[gray!60,dashed,line width=0.6pt] (0.134,1.317) -- (0.134,-1.317);
  \node[gray,scale=0.7,rotate=90,anchor=south] at (0.20,0) {отражение};
  \node[pt] at (-1,1.732) {}; \node[scale=0.85,anchor=south east] at (-1.03,1.76) {$P$};
  \node[pt] at (1,1) {};      \node[scale=0.85,anchor=south west] at (1.06,1.02) {$Q$};
  \node[pt] at (0.134,1.317) {}; \node[scale=0.85,anchor=south west] at (0.20,1.34) {$R$};
  \node[pt] at (0.134,-1.317) {};\node[scale=0.85,anchor=north west] at (0.20,-1.36) {$P+Q$};
  \node[scale=0.82] at (0.3,3.1) {Сложение точек: $P+Q$};
\end{scope}
% ---- справа: шесть точек на y^2 = x^3 + 1 ----
\begin{scope}[xshift=7.4cm,yshift=0cm,xscale=1.0,yscale=0.62]
  \draw[->,line width=0.5pt] (-2.1,0) -- (3.1,0) node[below left,scale=0.75] {$x$};
  \draw[->,line width=0.5pt] (0,-4.1) -- (0,4.4) node[below left,scale=0.75] {$y$};
  \draw (1,0.1) -- (1,-0.1) node[below,scale=0.7] {$1$};
  \draw (2,0.1) -- (2,-0.1) node[below,scale=0.7] {$2$};
  \draw[line width=0.9pt,domain=-1:2.45,smooth,samples=120,variable=\x]
        plot ({\x},{sqrt(max(\x*\x*\x+1,0))});
  \draw[line width=0.9pt,domain=-1:2.45,smooth,samples=120,variable=\x]
        plot ({\x},{-sqrt(max(\x*\x*\x+1,0))});
  \node[pt] at (2,3) {};   \node[scale=0.8,anchor=south west] at (2.08,3.1) {$P(2;3)$};
  \node[pt] at (0,1) {};   \node[scale=0.8,anchor=south west] at (0.12,1.05) {$2P(0;1)$};
  \node[pt] at (-1,0) {};  \node[scale=0.8,anchor=south east] at (-1.05,0.2) {$3P(-1;0)$};
  \node[pt] at (0,-1) {};  \node[scale=0.8,anchor=north west] at (0.12,-1.05) {$4P(0;-1)$};
  \node[pt] at (2,-3) {};  \node[scale=0.8,anchor=north west] at (2.08,-3.1) {$5P(2;-3)$};
  \node[gray,scale=0.78,align=left,anchor=north west] at (-2.05,4.15)
       {$6P=O$ --- бесконечно\\удалённая точка};
  \node[scale=0.82] at (0.6,5.0) {$y^2=x^3+1$: подгруппа $\langle P\rangle$};
\end{scope}
\end{tikzpicture}
\caption{Сложение точек на кубической кривой и пример конечной группы}
\label{fig:cubic}
\end{figure}

Коммутативность операции очевидна из симметрии определения, а обратным к точке
$P$ служит точка, симметричная $P$ относительно оси абсцисс. Проверка
ассоциативности нетривиальна: прямое вычисление требует разбора многих случаев,
а концептуальное доказательство использует теорему Безу или теорию дивизоров.
Поэтому ассоциативность здесь не доказывается и принимается со ссылкой
\cite[гл.~I, \S~4]{silverman}, \cite[гл.~7]{cassels}; см. также \cite{prasolov}.

Явные формулы сложения получаются из доказательства
предложения~\ref{prop:third}. Если $P(x_1,y_1)$ и $Q(x_2,y_2)$, $x_1\neq x_2$, то
угловой коэффициент секущей равен $\lambda=(y_2-y_1)/(x_2-x_1)$, третий корень
кубического уравнения равен $k^2-x_1-x_2$ при $k=\lambda$, и
\begin{equation}\label{eq:add}
  x_3=\lambda^2-x_1-x_2,\qquad y_3=\lambda(x_1-x_3)-y_1,\qquad P+Q=(x_3,y_3).
\end{equation}
Для удвоения вместо секущей берётся касательная. Дифференцируя уравнение кривой,
$2yy'=3x^2+a$, получаем, что для точки $P(x_0,y_0)$ с $y_0\neq0$ угловой
коэффициент касательной равен $\lambda=(3x_0^2+a)/(2y_0)$, а координаты точки
$2P$ суть
\begin{equation}\label{eq:double}
  x_2=\lambda^2-2x_0,\qquad y_2=\lambda(x_0-x_2)-y_0 .
\end{equation}
Если $y_0=0$, касательная вертикальна и $2P=O$.

\begin{exmp}
Кривая $y^2=x^3-2$ содержит рациональную точку $P(3,5)$. Здесь $a=0$ и
$\lambda=27/10$, откуда $2P=(129/100,\;-383/1000)$. Дальнейшее применение
операции порождает бесконечную последовательность рациональных точек с быстро
растущими знаменателями.
\end{exmp}

Таким образом, рациональные точки здесь не параметризуются: они порождаются
из конечного набора исходных точек повторением операции, и никакой рациональной
функции одного параметра, перечисляющей их все, не существует. Более того,
группа может оказаться и конечной~--- соответствующий пример разобран в
упражнении~\ref{ex:cubic}.

\section{Граница метода: теорема Морделла}

Окончательное описание множества рациональных точек кубики даёт следующая
теорема. Её доказательство использует метод спуска в сочетании с теорией высот
и выходит далеко за рамки настоящей работы.

\begin{thm}[Морделл, 1922]\label{thm:mordell}
Группа рациональных точек неособой кубической кривой, заданной над полем
рациональных чисел и имеющей хотя бы одну рациональную точку, конечно порождена
\cite[гл.~III]{silverman}, \cite[гл.~10]{cassels}.
\end{thm}

Иными словами, существует конечный набор рациональных точек, из которого все
остальные получаются описанной операцией. По теореме о строении конечно
порождённых абелевых групп группа рациональных точек изоморфна прямой сумме
конечной группы (группы кручения) и свободной группы некоторого ранга.

Сопоставим это со случаем коники. Для коники множество рациональных точек либо
пусто, либо параметризуется прямой; никакой групповой структуры для его описания
не требуется: все точки перечисляются одним рациональным параметром. Для
кубики параметризация невозможна, зато появляется группа, и содержательным
становится вопрос о её ранге~--- вопрос, далёкий от окончательного решения и по
сей день.

Наконец, для кривых рода не меньше двух положение меняется ещё раз: по теореме
Фальтингса, доказанной в 1983 году, множество их рациональных точек конечно
\cite[\S~6]{ostrik}. Итак, три разобранные ситуации отвечают трём
значениям рода кривой: для рода нуль множество рациональных точек пусто или
параметризуется прямой; для рода единица (при наличии рациональной точки) оно
образует конечно порождённую группу, которая может быть как бесконечной, так и
конечной (упражнение~\ref{ex:cubic}); для рода, большего единицы, оно всегда
конечно. Метод секущих
полностью решает задачу только в первом из этих случаев.

%=====================================================================
\chapter{ПРИМЕРЫ РЕШЕНИЯ ЗАДАЧ}
%=====================================================================

Для закрепления теоретического материала, изложенного в предыдущих главах,
рассмотрим решения нескольких типовых задач, связанных с методом секущих,
рациональными точками коник и их проективной интерпретацией.

\begin{exer}\label{ex:two}
\itshape Найти все рациональные точки окружности $x^2+y^2=2$, исходя из точки
$(1;1)$.
\end{exer}

\begin{proof}[Решение]
\normalfont
Точка $A(1,1)$ лежит на окружности, так как $1+1=2$. Проведём через неё прямую
с угловым коэффициентом $t$: $y=1+t(x-1)$. Подставим это выражение в уравнение
окружности:
\[
  x^2+\bigl[1+t(x-1)\bigr]^2=2 .
\]
Раскроем скобки и перенесём всё в левую часть:
\[
  x^2-1+2t(x-1)+t^2(x-1)^2=0 .
\]
Первое слагаемое разложим как $(x-1)(x+1)$ и вынесем общий множитель:
\[
  (x-1)\bigl[(x+1)+2t+t^2(x-1)\bigr]=0 .
\]
Корень $x=1$ отвечает исходной точке $A$. Второй корень находим из уравнения
$x(1+t^2)=t^2-2t-1$, откуда
\[
  x=\frac{t^2-2t-1}{t^2+1},\qquad
  y=1+t(x-1)=\frac{-t^2-2t+1}{t^2+1}.
\]
Проверим, что найденная точка действительно лежит на окружности. Числитель
суммы квадратов равен
\[
  (t^2-2t-1)^2+(t^2+2t-1)^2=2t^4+4t^2+2=2(t^2+1)^2,
\]
и после деления на $(t^2+1)^2$ получаем~2, что и требовалось.

При $t=0$ формулы дают точку $(-1;1)$, при $t=1$~--- точку $(-1;-1)$, при
$t=1/3$~--- точку $(-7/5;\,1/5)$, при $t=3$~--- точку $(1/5;\,-7/5)$.

Отметим два особых значения. Касательная к окружности в точке $A$ есть прямая
$x+y=2$ с угловым коэффициентом $t=-1$; при $t=-1$ формулы дают саму точку
$(1;1)$. В параметризацию не попадает вертикальная прямая $x=1$ (<<$t=\infty$>>):
она не касается окружности и пересекает её во второй точке $(1;-1)$, которую
формулы не дают ни при каком рациональном $t$.

\answer{рациональные точки окружности $x^2+y^2=2$~--- это точка $(1;-1)$ и точки
$x=\dfrac{t^2-2t-1}{t^2+1}$, $y=\dfrac{-t^2-2t+1}{t^2+1}$, где $t$ пробегает
множество рациональных чисел (значение $t=-1$ даёт саму точку $(1;1)$).}
\end{proof}

\begin{exer}\label{ex:seven}
\itshape Доказать, что окружность $x^2+y^2=7$ не содержит рациональных точек.
\end{exer}

\begin{proof}[Решение]
\normalfont
Рассуждаем так же, как в теореме~\ref{thm:three}. Предположим, что рациональная
точка существует. Записав её координаты с общим знаменателем, $x=X/Z$ и $y=Y/Z$,
где $X$, $Y$, $Z$ целые, $Z\neq0$ и $\gcdru(X,Y,Z)=1$, после
умножения на $Z^2$ получим
\[
  X^2+Y^2=7Z^2 .
\]
Выпишем все квадратичные вычеты по модулю~7. Возводя в квадрат остатки
$0,1,2,3,4,5,6$, получаем соответственно $0,1,4,2,2,4,1$. Значит, квадрат целого
числа сравним по модулю~7 с одним из чисел $0$, $1$, $2$, $4$.

Правая часть равенства сравнима с нулём по модулю~7, поэтому
$X^2+Y^2\equiv0\pmod 7$. Переберём все суммы пар из множества $\{0,1,2,4\}$ по
модулю~7:
\[
  0+0=0,\quad 0+1=1,\quad 0+2=2,\quad 0+4=4,\quad 1+1=2,
\]
\[
  1+2=3,\quad 1+4=5,\quad 2+2=4,\quad 2+4=6,\quad 4+4=1 .
\]
Нулю сравнима только сумма $0+0$.

Следовательно, $X^2\equiv0$ и $Y^2\equiv0 \pmod 7$, а поскольку число~7 простое,
отсюда 7 делит $X$ и 7 делит $Y$. Положим $X=7X_1$, $Y=7Y_1$. Подстановка даёт
$49X_1^2+49Y_1^2=7Z^2$, то есть $7(X_1^2+Y_1^2)=Z^2$. Значит, 7 делит $Z^2$, и,
поскольку 7~--- простое число, 7 делит $Z$. Но тогда $X$, $Y$ и $Z$ имеют общий
делитель~7, что противоречит условию $\gcdru(X,Y,Z)=1$.

\answer{рациональных точек на окружности $x^2+y^2=7$ нет. Отметим причину:
ключевым свойством оказалось то, что число $-1$ не является квадратичным вычетом
по модулю~7.}
\end{proof}

\begin{exer}\label{ex:pell}
\itshape Методом секущих найти три решения уравнения $x^2-3y^2=1$ в натуральных
числах.
\end{exer}

\begin{proof}[Решение]
\normalfont
Уравнение задаёт гиперболу, на которой лежит очевидная рациональная точка
$P_0(1,0)$. Проведём через неё прямую $y=t(x-1)$ и подставим в уравнение кривой:
\[
  x^2-3t^2(x-1)^2=1,\qquad\text{то есть}\qquad (x-1)(x+1)=3t^2(x-1)^2 .
\]
Сокращая на $x-1$, что законно при $x\neq1$, получаем $x+1=3t^2(x-1)$, откуда
$x(3t^2-1)=3t^2+1$ и
\[
  x=\frac{3t^2+1}{3t^2-1},\qquad y=\frac{2t}{3t^2-1}.
\]
Будем подставлять рациональные значения параметра, подбирая их так, чтобы
знаменатель $3t^2-1$ сокращался и точка оказывалась целочисленной.

При $t=1$ знаменатель равен~2, и мы получаем $x=4/2=2$, $y=2/2=1$. Проверка:
$4-3=1$.

При $t=2/3$ имеем $3t^2=4/3$, знаменатель равен $1/3$, и
$x=(4/3+1)/(1/3)=7$, $y=(4/3)/(1/3)=4$. Проверка: $49-48=1$.

При $t=3/5$ имеем $3t^2=27/25$, знаменатель равен $2/25$, и
$x=(52/25)/(2/25)=26$, $y=(6/5)/(2/25)=15$. Проверка: $676-675=1$.

\answer{$(2;1)$, $(7;4)$, $(26;15)$. Отметим, что параметризация даёт все
рациональные точки гиперболы, но выделение среди них целочисленных требует
подбора параметра; систематическое описание целых решений уравнения Пелля
методом секущих не получается.}
\end{proof}

\begin{exer}\label{ex:div3}
\itshape Доказать, что в любой примитивной пифагоровой тройке ровно один из
катетов делится на~3.
\end{exer}

\begin{proof}[Решение]
\normalfont
По теореме~\ref{thm:triples} всякая примитивная тройка имеет вид $a=p^2-q^2$,
$b=2pq$, $c=p^2+q^2$, где $\gcdru(p,q)=1$ и числа $p$, $q$ разной
чётности. Рассмотрим остатки $p$ и $q$ по модулю~3.

\emph{Случай первый: ни $p$, ни $q$ не делятся на~3.} Квадрат числа, не
делящегося на~3, сравним с~1 по модулю~3, поэтому $p^2\equiv q^2\equiv1\pmod 3$
и $a=p^2-q^2\equiv0\pmod 3$. Значит, катет $a$ делится на~3.

\emph{Случай второй: хотя бы одно из чисел $p$, $q$ делится на~3.} Тогда
произведение $pq$ делится на~3, а с ним и катет $b=2pq$.

Итак, хотя бы один из катетов всегда делится на~3. Остаётся показать, что оба
одновременно делиться не могут. Предположим противное: пусть 3 делит и
$a=p^2-q^2$, и $b=2pq$. Из делимости $2pq$ на~3 и простоты числа~3 следует, что
3 делит $p$ или 3 делит $q$. Пусть для определённости 3 делит $p$. Тогда
$p^2\equiv0\pmod 3$, и из делимости $p^2-q^2$ на~3 получаем, что 3 делит $q^2$,
а значит и $q$. Но тогда $p$ и $q$ имеют общий делитель~3, что противоречит
условию $\gcdru(p,q)=1$. Случай, когда 3 делит $q$, разбирается
симметрично.

\answer{ровно один катет примитивной тройки кратен трём. Например, в тройке
$(3,4,5)$ это первый катет, в тройке $(5,12,13)$~--- второй.}
\end{proof}

\begin{exer}\label{ex:infty}
\itshape Найти точки пересечения с бесконечно удалённой прямой для эллипса
$x^2+2y^2=1$, гиперболы $x^2-2y^2=1$ и параболы $y=x^2$.
\end{exer}

\begin{proof}[Решение]
\normalfont
Перейдём к однородным координатам, положив $x=x_1/x_0$ и $y=x_2/x_0$. Бесконечно
удалённая прямая задаётся уравнением $x_0=0$, поэтому в каждом случае нужно
записать уравнение кривой в однородных координатах и положить $x_0=0$.

\emph{Эллипс.} Умножая на $x_0^2$, получаем $x_1^2+2x_2^2=x_0^2$. При $x_0=0$
остаётся $x_1^2+2x_2^2=0$. Над полем вещественных чисел оба слагаемых
неотрицательны, поэтому равенство возможно лишь при $x_1=x_2=0$, но нулевая
тройка не задаёт точку проективной плоскости. Вещественных бесконечно удалённых
точек нет.

\emph{Гипербола.} Аналогично получаем $x_1^2-2x_2^2=x_0^2$, и при $x_0=0$
остаётся $x_1^2=2x_2^2$. Отсюда $x_1=\pm\sqrt2\,x_2$, что даёт две различные
точки: $(0:\sqrt2:1)$ и $(0:-\sqrt2:1)$. Обе вещественны, но ни одна не
рациональна, так как число $\sqrt2$ иррационально. Это замечание существенно:
при описании рациональных точек гиперболы бесконечно удалённые точки не
возникают.

\emph{Парабола.} Уравнение $y=x^2$ в однородных координатах принимает вид
$x_2x_0=x_1^2$. При $x_0=0$ получаем $x_1^2=0$, откуда $x_1=0$ и остаётся
единственная точка $(0:0:1)$. Корень $x_1=0$ двукратен, поэтому бесконечно
удалённая прямая касается параболы в этой точке.

\answer{эллипс не имеет вещественных бесконечно удалённых точек, гипербола имеет
две (иррациональные), парабола касается бесконечно удалённой прямой в одной
точке (см. рис.~\ref{fig:infinity}). В этом и состоит проективное различие трёх типов
коник: как проективные кривые они не различаются, различается лишь их положение
относительно выбранной прямой.}
\end{proof}

\begin{exer}\label{ex:cubic}
\itshape На кривой $y^2=x^3+1$ дана точка $P(2;3)$. Последовательно вычислить
$2P$, $3P$ и $4P$ и определить порядок точки $P$ в группе рациональных точек.
\end{exer}

\begin{proof}[Решение]
\normalfont
Точка $P$ лежит на кривой: $3^2=9$ и $2^3+1=9$. Воспользуемся формулами
\eqref{eq:add} и \eqref{eq:double} (здесь $a=0$); итог вычислений изображён на рис.~\ref{fig:cubic} справа.

\emph{Вычислим $2P$.} Угловой коэффициент касательной в точке $P$ равен
$\lambda=(3x^2+a)/(2y)=3\cdot4/(2\cdot3)=2$. Тогда
\[
  x(2P)=\lambda^2-2x=4-4=0,\qquad y(2P)=\lambda(x-x(2P))-y=2(2-0)-3=1,
\]
то есть $2P=(0;1)$. Проверка: $1^2=1$ и $0^3+1=1$.

\emph{Вычислим $3P=P+2P$.} Проводим прямую через точки $(2;3)$ и $(0;1)$; её
угловой коэффициент $\lambda=(1-3)/(0-2)=1$. Тогда
\[
  x(3P)=\lambda^2-2-0=1-2=-1,\qquad y(3P)=\lambda\bigl(2-(-1)\bigr)-3=3-3=0,
\]
то есть $3P=(-1;0)$. Проверка: $0^2=0$ и $(-1)^3+1=0$.

\emph{Вычислим $4P=2P+2P$}, то есть удвоим точку $(0;1)$. Здесь
$\lambda=3\cdot0^2/(2\cdot1)=0$, откуда
\[
  x(4P)=0^2-2\cdot0=0,\qquad y(4P)=0\cdot(0-0)-1=-1,
\]
то есть $4P=(0;-1)$. Но это точка, симметричная $2P$ относительно оси абсцисс,
то есть $4P=-2P$. Значит, $6P=O$, где $O$~--- бесконечно удалённая точка.

Осталось убедиться, что меньший порядок невозможен. Точки $P$, $2P$ и $3P$
отличны от $O$, поскольку все они имеют конечные координаты; $4P=-2P\neq O$ и
$5P=-P\neq O$. Следовательно, наименьшее натуральное $n$, для которого $nP=O$,
равно шести.

\answer{$2P=(0;1)$, $3P=(-1;0)$, $4P=(0;-1)$; порядок точки $P$ равен~6. Эта
задача показывает то, что для коник невозможно по следствию~\ref{cor:dich}: у
кубической кривой множество рациональных точек может быть конечным и ненулевым.
Подгруппа, порождённая точкой $P$, состоит из шести точек: $(2;3)$, $(0;1)$,
$(-1;0)$, $(0;-1)$, $(2;-3)$ и $O$. То, что других рациональных точек на этой
кривой нет, следует из теоремы Нагелля~--- Лутц и равенства нулю ранга
\cite[гл.~II]{silverman}; здесь этот факт не доказывается.}
\end{proof}

%=====================================================================
\chapter*{ЗАКЛЮЧЕНИЕ}
\addcontentsline{toc}{chapter}{ЗАКЛЮЧЕНИЕ}
%=====================================================================

В работе изложена рациональная параметризация конических сечений и прослежены её
основания и границы.

Разбор начат с классической задачи о пифагоровых тройках, которая сведена к
описанию рациональных точек единичной окружности. Выделена и доказана лемма о
втором корне квадратного уравнения~--- элементарное утверждение, на котором
держится вся дальнейшая конструкция. Метод секущих даёт полное решение задачи о
тройках и приводит к классической формуле для примитивных троек.

Показано, что конструкция не связана с окружностью и переносится на произвольную
неособую конику с рациональными коэффициентами: наличие хотя бы одной
рациональной точки влечёт бесконечность их множества, причём все они получаются
секущими. Отсюда следует, что для коники возможны ровно две ситуации~--- пустое
множество рациональных точек либо бесконечное; промежуточного случая нет. Вопрос
о том, какая из ситуаций реализуется, решается принципом Хассе~--- Минковского;
приведён элементарный пример коники $x^2+y^2=3$, не имеющей рациональных точек.

Проективная переформулировка показывает, почему метод работает: коника с
отмеченной рациональной точкой есть проективный образ пучка прямых с центром в
этой точке, что является частным случаем теоремы Штейнера. В проективной
постановке исчезают исключительные случаи аффинной формулировки, а итоговое
утверждение принимает окончательный вид: неособая коника с рациональной точкой
рационально эквивалентна проективной прямой.

Показано, что при переходе к кривым третьего порядка метод перестаёт быть
исчерпывающим: секущие дают не параметризацию, а групповую операцию, а описание
множества рациональных точек даётся теоремой Морделла о конечной порождённости.
Тем самым очерчена область применимости метода~--- кривые рода нуль.

В пятой главе разобраны шесть задач, охватывающих материал всех теоретических
глав: применение метода секущих к конике с нестандартной отмеченной точкой,
доказательство отсутствия рациональных точек методом сравнений, получение
решений уравнения Пелля, свойство делимости катетов примитивной тройки,
вычисление бесконечно удалённых точек коник и построение конечной группы точек
на кубической кривой. Последняя задача показывает, что для
кубики, в отличие от коники (следствие~\ref{cor:dich}), конечное ненулевое
число рациональных точек возможно.

Все утверждения, допускающие элементарное обоснование, доказаны в работе
самостоятельно; результаты, требующие аппарата за рамками настоящей работы,
приведены с указанием источников. Поставленные во введении задачи решены, цель
работы достигнута.

Естественным продолжением является подробное изучение группового закона на
эллиптических кривых, а также методическая разработка материала первой главы для
занятий со школьниками, поскольку он требует лишь знания уравнения окружности и
умения решать квадратное уравнение.

%=====================================================================
%                        СПИСОК ЛИТЕРАТУРЫ
%=====================================================================
\renewcommand{\bibname}{СПИСОК ЛИТЕРАТУРЫ}
\begin{thebibliography}{99}
\addcontentsline{toc}{chapter}{СПИСОК ЛИТЕРАТУРЫ}

\bibitem{atanasyan} Атанасян Л.~С., Базылев В.~Т. Геометрия. Часть~II. ---
  Москва: Просвещение, 1987. --- 368~с.

\bibitem{borevich} Боревич З.~И., Шафаревич И.~Р. Теория чисел. --- 3-е изд. ---
  Москва: Наука, 1985. --- 504~с.

\bibitem{efimov} Ефимов Н.~В. Высшая геометрия. --- Москва: Наука, 1971. ---
  584~с.

\bibitem{coxeter} Кокстер Г.~С.~М. Действительная проективная плоскость. ---
  Москва: Физматгиз, 1959. --- 280~с.

\bibitem{ostrik} Острик В.~В., Цфасман М.~А. Алгебраическая геометрия и теория
  чисел: рациональные и эллиптические кривые. --- Москва: МЦНМО, 2001. --- 48~с.

\bibitem{prasolov} Прасолов В.~В., Соловьёв Ю.~П. Эллиптические функции и
  алгебраические уравнения. --- Москва: Факториал, 1997. --- 288~с.

\bibitem{serre} Серр Ж.-П. Курс арифметики. --- Москва: Мир, 1972. --- 184~с.

\bibitem{silverman} Сильверман Дж., Тейт Дж. Рациональные точки на эллиптических
  кривых. --- Москва: Мир, 2000. --- 448~с.

\bibitem{vinogradov} Виноградов И.~М. Основы теории чисел. --- 9-е изд. ---
  Москва: Наука, 1981. --- 176~с.

\bibitem{hartshorne} Хартсхорн Р. Основы проективной геометрии. --- Москва: Мир,
  1970. --- 160~с.

\bibitem{cassels} Cassels J.~W.~S. Lectures on Elliptic Curves. --- Cambridge:
  Cambridge University Press, 1991. --- 137~p.

\bibitem{hatcher} Hatcher A. Topology of Numbers. --- Providence: American
  Mathematical Society, 2022. --- 341~p.

\end{thebibliography}

%=====================================================================
\chapter*{ПРИЛОЖЕНИЕ}
\addcontentsline{toc}{chapter}{ПРИЛОЖЕНИЕ}
%=====================================================================

В рамках курсовой работы разработана интерактивная модель, иллюстрирующая метод
секущих и его применения. Модель выполнена в виде веб-страницы и не требует
установки программ: достаточно открыть её по ссылке с компьютера или телефона.

\textbf{Адрес модели:} \url{https://mvbulgakova.github.io/kursovaya_mat_2026/}

\textbf{Исходный код:} \url{https://github.com/mvbulgakova/kursovaya_mat_2026}

Модель состоит из четырёх разделов.

\textbf{Секущие.} Через отмеченную рациональную точку коники проводится секущая
с угловым коэффициентом $t=p/q$ и вычисляется вторая точка пересечения. Доступны
окружности $x^2+y^2=1$ и $x^2+y^2=2$, гипербола $x^2-2y^2=1$ и эллипс
$x^2+xy+y^2=1$; для окружности сразу выводится соответствующая пифагорова тройка,
для эллипса~--- треугольник с углом $120^\circ$ или $60^\circ$. В режиме анимации
секущие перебирают дроби $p/q$ по возрастанию сложности, и на экране видно, как
рациональные точки постепенно заполняют кривую. Для окружности $x^2+y^2=3$
секущие провести не из чего; вместо них модель выполняет перебор, не находящий
ни одной рациональной точки (теорема~\ref{thm:three}).

\textbf{Целые треугольники.} Выводятся все прямоугольные треугольники, треугольники
с углом $120^\circ$ или $60^\circ$ и героновы треугольники (с целой площадью),
стороны которых не превосходят заданного числа. Для выбранного треугольника
показан чертёж и соответствующая ему точка на конике вместе с секущей,
которая её даёт. Для героновых треугольников показана высота, разбивающая
треугольник на два прямоугольных с рациональными сторонами.

\textbf{Задачи для урока.} Генератор самостоятельных работ по вариантам: нахождение
гипотенузы и катета, расстояния между точками на координатной плоскости,
третьей стороны по теореме косинусов при угле $120^\circ$ или $60^\circ$, площади
по формуле Герона, целых точек на окружности. Числа в задачах берутся из
пифагоровых и эйзенштейновых троек, поэтому все ответы получаются целыми.
Варианты можно распечатать вместе с листом ответов.

\textbf{Кубики.} Вычисляются кратные $nP$ точки на кривых $y^2=x^3+1$ (точка $P(2;3)$
порядка шесть) и $y^2=x^3-2$ (точка $P(3;5)$ бесконечного порядка) при $n$ до~200.
График показывает, что число цифр в знаменателе растёт примерно как $n^2$: у
точки $200P$ знаменатель абсциссы содержит более 23\,000 цифр. Чтобы такие
вычисления занимали доли секунды, координаты $nP$ находятся не последовательным
сложением, а через значения полиномов деления:
$x(nP)=x-\psi_{n-1}\psi_{n+1}/\psi_n^2$, где целые числа $\psi_n$ задаются
рекуррентными формулами.

Все вычисления в модели выполняются в точной арифметике рациональных чисел над
длинными целыми, без округления. Корректность проверена отдельной программой
(файл \texttt{model/test/verify.js}): для четырёх коник, всех 183 несократимых
значений параметра с $|p|,q\leqslant12$ и вертикальной прямой (всего 736 точек)
подстановка координат в уравнение кривой даёт точный нуль; кроме того,
пересчитаны все числовые примеры основного текста.

Вторая точка пересечения вычисляется по общей формуле, пригодной для
произвольной коники \eqref{eq:conic} с отмеченной точкой $(x_0;y_0)$. Подстановка
$y=y_0+t(x-x_0)$ и замена $u=x-x_0$ приводят уравнение к виду
$u\cdot\bigl(A(t)\,u+B(t)\bigr)=0$, где
\[
  A(t)=a+bt+ct^2,\qquad
  B(t)=2a\,x_0+b(x_0t+y_0)+2c\,y_0t+d+et,
\]
причём свободный член обращается в нуль, поскольку точка $(x_0;y_0)$ лежит на
конике. Отсюда $x_1=x_0-B/A$ и $y_1=y_0-t\cdot B/A$. Случай $A(t)=0$
обрабатывается отдельно: вторая точка пересечения бесконечно удалена, о чём
модель сообщает.

Чертежи, помещённые в основной текст работы (рисунки \ref{fig:secant},
\ref{fig:infinity} и \ref{fig:cubic}), построены средствами пакета Ti\emph{k}Z
непосредственно в исходном коде документа и не требуют внешних графических
файлов.

\end{document}
